\section{Introduction}

\label{sec:introduction}
% ══════════════════════════════════════════════════════════════════════════════

The voluntary carbon market (VCM) reached \$2 billion in 2021
\citep{foresttrends2022}, with nature-based solutions (NBS) representing
the fastest-growing segment. NBS projects---including avoided
deforestation (REDD+), reforestation, mangrove restoration, and peatland
conservation---generate carbon credits by sequestering atmospheric CO$_2$
or preventing its release from terrestrial carbon stocks.

However, the scientific credibility of NBS carbon credits has been
severely challenged. \citet{west2023action} found that 94\% of Verra's
rainforest offset credits did not represent genuine emissions reductions.
\citet{guizar2022carbon} showed that certified REDD+ projects in the
Brazilian Amazon overestimated avoided deforestation by 1.8x on average.
These findings have eroded buyer confidence and threaten the viability of
carbon markets as a climate mitigation tool.

The fundamental problem is MRV infrastructure. Current practice relies
on:

\begin{enumerate}[label=(\roman*)]
  \item \textbf{Infrequent verification}: Third-party audits occur every
    2--5 years, leaving extended periods without monitoring
    \citep{schneider2019double}.
  \item \textbf{Self-reported baselines}: Project developers construct
    their own counterfactual deforestation scenarios, creating systematic
    upward bias \citep{west2023action}.
  \item \textbf{Manual field sampling}: Biomass estimation relies on
    small plot samples extrapolated to vast areas, introducing large
    sampling errors \citep{chave2014improved}.
  \item \textbf{Opaque methodologies}: Verification reports are often
    proprietary, preventing independent scrutiny \citep{badgley2022}.
  \item \textbf{Delayed detection}: Reversals (fire, illegal logging,
    drought mortality) may go undetected for years between audit cycles.
\end{enumerate}

\subsection{Digital MRV}

Advances in satellite remote sensing, machine learning, and distributed
ledger technology enable a fundamentally different approach to MRV---one
that is continuous, transparent, and independently verifiable. Digital
MRV (dMRV) replaces periodic manual audits with automated monitoring
pipelines that process satellite data streams in near-real-time
\citep{fuss2020moving}.

\subsection{Contributions}

We present ZooMRV, an end-to-end digital MRV platform with the following
contributions:

\begin{enumerate}
  \item \textbf{CarbonNet}: A deep learning model for above-ground
    biomass (AGB) estimation from multi-spectral satellite imagery and
    LiDAR, achieving RMSE of 18.4 Mg/ha on independent validation data.

  \item \textbf{Dynamic baselining}: An algorithmic approach to
    counterfactual construction using synthetic controls
    \citep{abadie2010synthetic} that eliminates self-reported baseline
    bias.

  \item \textbf{Continuous monitoring}: Near-real-time change detection
    with 12-day median latency for reversal events, compared to 14
    months for traditional audit cycles.

  \item \textbf{On-chain attestation}: A blockchain-based credit
    registry where each credit is linked to verifiable satellite data,
    spatial boundaries, and algorithmic verification reports.

  \item \textbf{Open validation}: All models, data, and methodologies
    are open-source, enabling independent verification and continuous
    improvement by the scientific community.
\end{enumerate}


% ══════════════════════════════════════════════════════════════════════════════
